%% kapa: exam=1 round=37 year=2026 period=2 form=A subject=accounting subject_ko=회계학 seq=40 number=80
%% source: Q-Net 감정평가사 기출문제 (cst00309) artlSeq=5258924 file=raw/1008/5258924/2026년 제37회 감정평가사 1차 2교시.pdf
%% conversion: kordoc 4.24.51 -> tools/gen1.py (auto); review=proofread-sonnet; flags=none
\begin{question}{80}
(주)감평은 단일제품을 생산ㆍ판매하고 있으며, 시장수요예측에 근거하여 당기에 21,000단위의 제품을 판매할 계획이다. 원재료 X는 공정 초에 전량 투입(제품 단위당 2kg 투입)되고 원재료 Y는 공정의 50 \%시점에 전량 투입(제품 단위당 3kg 투입)되며, 가공원가는 공정 전반에 걸쳐 균등하게 발생한다. 당기 초 실제재고와 당기 말 목표재고는 다음과 같다.
\par
{\small\centering
\begin{tabular}{@{}lcc@{}}\hline
 & 기초실제재고 & 기말목표재고 \\\cline{2-3}
원재료 X & 1,800kg & 2,400kg \\
원재료 Y & 3,000kg & 4,500kg \\
재 공 품 & 800단위(가공원가 완성도 25 \%) & 800단위(가공원가 완성도 75 \%) \\
제 \quad 품 & 700단위 & 400단위 \\\hline
\end{tabular}\par}
(주)감평이 선입선출법에 의해 재고자산을 평가할 때, 당기에 구입해야 하는 각 원재료의 수량은? (단, 제조과정에서 공손 및 감손 등으로 인한 물량손실은 발생하지 않는다.) \par
\hspace*{2.4em}\makebox[4.5cm][l]{\uline{원재료 X}}\uline{원재료 Y} \par
\begin{choices}
  \item \makebox[4.5cm][l]{42,000kg,}63,600kg
  \item \makebox[4.5cm][l]{42,000kg,}64,800kg
  \item \makebox[4.5cm][l]{42,000kg,}66,000kg
  \item \makebox[4.5cm][l]{42,800kg,}64,800kg
  \item \makebox[4.5cm][l]{42,800kg,}66,000kg
\end{choices}
\end{question}
\answer{3}
